逐步尝试后发现，TF-IDF 相对 BoW 使 SVM、LR、MLP 的最佳 F1 分别提高 3.22、2.02、0.27 个百分点，NB 也提高 0.15 个百分点。当前 TF-IDF 训练矩阵有 15,000 维，但非零元素只占 0.80\%；每篇文本通常只激活少量词项，类别又可由一组区分词形成近似线性边界。SVM 直接寻找大间隔线性边界，LR 则用正则化权重累积各词项对类别的贡献，因此都适合这种高维稀疏输入。NB 的条件独立假设虽然不完全成立，但新闻主题常由若干高区分度词共同指示；多项式 NB 可以直接累积这些非负词项证据，而较均衡的类别分布也减弱了先验偏置，因此以很小的计算量取得了 93.51\% 的 F1。MLP 的 94.22\% 最高，但拟合约需 99 秒；LR 用约 5 秒达到 93.48\%。

参数搜索也改变了比较结论。TF-IDF＋SVM 的 $C$ 从 0.1 增至 1 时，F1 从 84.72\% 升至 93.23\%，继续增至 10 反而略降；LR 的 $C=10$ 与 100 分别为 93.47\% 和 93.48\%，补查 1,000 后降至 93.35\%。平均 Word2Vec＋LR 最高为 87.45\%，词向量构造另耗时约 213 秒。这只说明当前小语料上“从头训练＋均值聚合”的效果，不代表预训练词向量或 BERT。
